\section{Sequencer}\label{sequencer}

The LXR has a sequencer for playing and programming drum patterns. The
basic functionality is similar to classic TR-x0x machines, with
additional features.

\subsection{Tracks}\label{tracks}

There are 7 sequencer tracks --- one per voice, plus an additional track
to control the open hihat of voice 6.

\subsubsection{Select Tracks}\label{select-tracks}

The active track is changed with the 7 voice buttons. A lit LED next to
the button shows the selected track.

\subsubsection{Mute Tracks}\label{mute-tracks}

Hold the shift button and press a voice button to mute/unmute that
track. The LEDs of unmuted voices light up.

\tip{In performance mode the function of the voice buttons is inverted — they mute/unmute
without needing to hold shift. This allows single-handed operation.}

\subsubsection{Unmute All Tracks}\label{unmute-all-tracks}

To unmute all muted tracks simultaneously, go to performance mode and
hold shift. The voice buttons then unmute all tracks up to and including
the pressed button.

Example:

\begin{itemize}
\tightlist
\item
  Shift + Voice Button 4 = tracks 1, 2, 3 and 4 will be unmuted.
\item
  Shift + Voice Button 7 = all tracks are unmuted.
\end{itemize}

\subsection{Steps}\label{steps}

There are 16 steps per track. Steps are set or removed using the 16
sequencer buttons.

\subsubsection{Sub Steps}\label{sub-steps}

Each of the 16 main steps consists of 8 sub steps, giving a total of 128
steps per pattern.

\begin{figure}[H]
  \centering
  \includegraphics[width=\linewidth]{substep_diagram.png}
  \caption{Sub steps shown within the step buttons.}
\end{figure}

Sub steps define a step pattern played when the corresponding main step
is active. They can be used to create flam or roll sounds on a
particular step, or to place a note between the \(\frac{1}{16}\)th notes
of the main steps.

If the main step is inactive, none of its sub steps will be played. Each
sub step can have its own pitch, volume, probability, and automation
settings.

\subsection{Pattern}\label{pattern}

A set of all 7 tracks is called a pattern. Different patterns can be
chained to play one after another --- see the Performance Mode section
for details.

\subsection{Pattern Set}\label{pattern-set}

A pattern set is a group of 8 patterns saved together on the SD card and
directly accessible in performance mode.

\subsection{Basic Pattern Editing}\label{basic-pattern-editing}

Basic pattern editing is done in voice edit mode. Select the track with
the voice buttons. The step data is displayed on the 16 sequencer button
LEDs. Use the sequencer buttons to set and remove steps.

\subsection{Step Parameters}\label{step-parameters}

Each step has a set of parameters. To edit them in voice edit mode, hold
the shift button. The display then shows the step parameter page and the
LED of the selected step starts flashing.

\subsubsection{Available Step
Parameters}\label{available-step-parameters}

\begin{paramtable}
vel & Step Velocity    & Note-on velocity (0–127). Can be used as a mod source on the
                         voice modulation page. \\
nte & Step Note        & MIDI note number to play (--63 to +63). Default 0 = MIDI note 63. \\
prb & Step Probability & Chance of the step being played. 0 = never (0\%), 63 = 50\%,
                         127 = always (100\%). \\
\end{paramtable}

\subsubsection{Step Probability}\label{step-probability}

The step probability value sets the chance of a step being played, from
0 to 127 (0--100\%).

On each main step, the sequencer generates a new random number between 0
and 127. If the random number is less than or equal to the step's
probability value, the step plays.

\tip{The random number is generated once per main step and shared across all 8 of its sub steps.
This allows interesting sub-step cluster randomisation — the 1st sub step always plays, the 5th
plays with 50\% chance, and the 3rd and 7th play with 25\% chance.}

\subsection{Automation}\label{automation}

Parameter values can be recorded to the sequencer steps. Each step can
store 2 parameter changes. They only affect that single step and are
ignored if the step is inactive. The modulation targets can be different
on every step.

\subsubsection{Choosing the Active Automation
Track}\label{choosing-the-active-automation-track}

There are 2 automation tracks per voice. Select the active track before
recording.

\begin{itemize}
\tightlist
\item
  Hold shift and press the record button to open the recording options
  menu.
\item
  The \param{trk} parameter selects the active automation track.
\end{itemize}

\subsubsection{How to Assign Parameter
Automations}\label{how-to-assign-parameter-automations}

\textbf{Realtime recording:}

\begin{enumerate}
\def\labelenumi{\arabic{enumi}.}
\tightlist
\item
  Press the record button. The record LED lights up.
\item
  While the pattern is playing, tweak sound parameters with the knobs.
  Changes are recorded to the played steps.
\end{enumerate}

\attention{Activate the quantisation grid before recording automations. Without quantisation,
data is recorded to all 128 sub steps, and most of it will land on inactive sub steps if you are
only using the 16 main steps.}

\textbf{Lock parameter values to a single step:}

\begin{enumerate}
\def\labelenumi{\arabic{enumi}.}
\tightlist
\item
  In voice edit mode, hold down the step button.
\item
  The step LED starts blinking.
\item
  Use one of the 4 knobs to assign a parameter value to that step.
\item
  Release the step button.
\end{enumerate}

\textbf{Edit automation manually via the step edit menu:}

\begin{enumerate}
\def\labelenumi{\arabic{enumi}.}
\tightlist
\item
  Go to the step parameter menu (voice mode: shift + step button).
\item
  Use the encoder to scroll to page 2. You will see:
\end{enumerate}

\begin{paramtable}
p1d & Parameter 1 Destination & Select the target voice parameter from a list. \\
p1v & Parameter 1 Value       & Set the parameter to this value on this step. \\
p2d & Parameter 2 Destination & Select the target voice parameter from a list. \\
p2v & Parameter 2 Value       & Set the parameter to this value on this step. \\
\end{paramtable}

Manual editing is the only way to assign parameter automations from one
track to another. Edit the target parameter in detail mode (push
encoder) to see full names and destination voice.

\subsubsection{Clear Automation Data}\label{clear-automation-data}

\begin{itemize}
\tightlist
\item
  Hold shift and press the copy/clear button.
\item
  Use the encoder to select \texttt{Clear\ {[}autom.1/2{]}?}
\item
  Press copy/clear again to confirm, or release shift to abort.
\end{itemize}

\subsection{Recording Options Menu}\label{recording-options-menu}

Hold shift and press the record button to open the recording options
menu:

\oled{Trk \quad qnt}{0 \hspace{2.5em} 16}

\subsubsection{Automation Track}\label{automation-track}

The \param{trk} parameter selects which of the 2 automation tracks is
active. Parameter automations are recorded to the active track.

\subsubsection{Quantisation}\label{quantisation}

Selects the quantisation grid. All recorded notes and automation data
will be quantised to the selected grid. Quantisation cannot be applied
retroactively --- it must be on before recording.

{\def\LTcaptype{none} % do not increment counter
\begin{longtable}[]{@{}ll@{}}
\toprule\noalign{}
Parameter value & Description \\
\midrule\noalign{}
\endhead
\bottomrule\noalign{}
\endlastfoot
off & No quantisation \\
8 & Quantise to 1/8 notes \\
16 & Quantise to 1/16 notes \\
32 & Quantise to 1/32 notes \\
64 & Quantise to 1/64 notes \\
\end{longtable}
}

\subsection{Copy Sequencer Data}\label{copy-sequencer-data}

It is possible to copy sequencer tracks and patterns to another
position.

\begin{enumerate}
\def\labelenumi{\arabic{enumi}.}
\tightlist
\item
  Press and hold the Copy button. Its LED starts flashing to indicate
  copy mode is active.
\item
  Select the source data by pressing a voice button (track copy) or a
  pattern button (pattern copy). The source LED starts flashing.
\item
  Press the target voice/pattern button. The data is copied immediately,
  overwriting the destination.
\item
  To abort, release the copy button before selecting a target.
\end{enumerate}

\subsection{Clear Sequencer Data}\label{clear-sequencer-data}

Sequencer data can be cleared separately for track, pattern, or
automation data.

\begin{enumerate}
\def\labelenumi{\arabic{enumi}.}
\tightlist
\item
  Press Shift + Copy/Clear.
\item
  The display shows \texttt{Clear\ {[}track{]}?}
\item
  Use the encoder to select what to delete: \texttt{{[}track{]}},
  \texttt{{[}pattern{]}}, \texttt{{[}autom.1{]}}, or
  \texttt{{[}autom.2{]}}.
\item
  Press the clear button again to confirm, or release shift to abort.
\end{enumerate}
